\documentclass[aspectratio=169,11pt]{beamer}
\usepackage{../../../shared/theme/esmad_beamer_theme}
\usepackage{amsmath,amssymb}
\usepackage{booktabs}

\title[Missing Data and Selection Bias]{Missing Data, Attrition, and Selection Bias in Causal Inference}
\subtitle{Identification When Outcomes or Units Disappear}
\date{\today}

\newcommand{\E}{\mathbb{E}}

\begin{document}

\begin{frame}
\titlepage
\end{frame}

\begin{frame}{Learning Goals}
\begin{itemize}
\item distinguish MCAR, MAR, and MNAR mechanisms;
\item explain how attrition can destroy randomized comparisons;
\item recognize complete-case analysis as an identification choice;
\item use inverse-probability-of-observation weighting and imputation appropriately;
\item understand Lee bounds and partial-identification logic;
\item design sensitivity analyses for missing-not-at-random outcomes.
\end{itemize}
\end{frame}

\begin{frame}{Outline}
\tableofcontents
\end{frame}

\section{Why Missingness Is Causal}

\begin{frame}{Observed Outcomes Are Selected Outcomes}
Let
\[
R_i=
\begin{cases}
1,&Y_i\text{ observed},\\
0,&Y_i\text{ missing}.
\end{cases}
\]

We observe \(R_iY_i\), not always \(Y_i\).

\begin{alertblock}{Key point}
The mechanism determining \(R_i\) can induce selection bias even when treatment assignment itself was randomized.
\end{alertblock}
\end{frame}

\begin{frame}{Attrition Can Break Randomization}
Suppose treatment is randomized, but
\[
P(R=1\mid D=1) \neq P(R=1\mid D=0).
\]

The observed treated and control samples may no longer be comparable if attrition depends on potential outcomes or post-treatment variables.
\end{frame}

\section{MCAR, MAR, and MNAR}

\begin{frame}{Missing Completely at Random}
MCAR means, informally,
\[
R \perp (Y,D,X).
\]

Missingness is unrelated to observed or unobserved variables relevant to the analysis.
\end{frame}

\begin{frame}{Missing at Random}
MAR means
\[
R \perp Y_{mis} \mid Y_{obs},D,X.
\]

After conditioning on observed information, missingness does not depend on the missing values themselves.
\end{frame}

\begin{frame}{Missing Not at Random}
MNAR means missingness still depends on the unobserved outcome after conditioning on observed information.

Example: participants with especially poor outcomes are more likely to disappear from follow-up.
\end{frame}

\section{Complete-Case Analysis}

\begin{frame}{What Complete-Case Analysis Assumes}
Analyzing only \(R=1\) changes the target population.

A complete-case contrast estimates
\[
\E[Y\mid D=1,R=1]-\E[Y\mid D=0,R=1].
\]

This equals the causal effect of interest only under strong selection assumptions.
\end{frame}

\begin{frame}{Selection as Conditioning on a Collider}
If both treatment and outcome affect observation,
\[
D \rightarrow R \leftarrow Y,
\]
conditioning on \(R=1\) can induce association between treatment and outcome.

\begin{alertblock}{Collider bias}
Observed cases can become a selected subpopulation created by the analysis itself.
\end{alertblock}
\end{frame}

\section{Inverse-Probability-of-Observation Weighting}

\begin{frame}{Observation Probabilities}
Estimate
\[
\pi(X,D)=P(R=1\mid X,D).
\]

Observed units receive weights
\[
w_i=\frac{1}{\widehat\pi(X_i,D_i)}.
\]
\end{frame}

\begin{frame}{Weighted Estimation}
Under MAR-type assumptions,
\[
\E\left[
\frac{R_iY_i}{\pi(X_i,D_i)}
\right]
\]
can recover population expectations.
\end{frame}

\begin{frame}{Weight Instability}
If \(\pi(X,D)\) is near zero, weights explode.

Consequences:
\begin{itemize}
\item high variance;
\item sensitivity to model misspecification;
\item effective sample-size collapse.
\end{itemize}
\end{frame}

\section{Imputation}

\begin{frame}{Outcome Imputation}
Under MAR, estimate
\[
m_d(X)=\E[Y\mid D=d,X,R=1]
\]
and predict missing outcomes.

Validity depends on the MAR assumption, outcome modeling, and overlap.
\end{frame}

\begin{frame}{Multiple Imputation}
For causal work:
\begin{itemize}
\item include treatment;
\item include strong outcome predictors;
\item include variables predicting missingness;
\item respect time ordering;
\item avoid post-treatment variables that change the estimand.
\end{itemize}
\end{frame}

\section{Doubly Robust Missingness Adjustment}

\begin{frame}{Outcome Model Plus Observation Model}
One can combine:
\begin{itemize}
\item a model for \(P(R=1\mid X,D)\);
\item a model for \(\E[Y\mid X,D,R=1]\).
\end{itemize}

Augmented estimators can be consistent if one nuisance model is correctly specified under suitable identification assumptions.
\end{frame}

\section{Lee Bounds}

\begin{frame}{Monotone Selection}
In randomized experiments with differential attrition, Lee bounds use a monotonicity assumption such as treatment weakly increasing observation probability.

The more-retained group is trimmed to match the retention rate of the other group.
\end{frame}

\begin{frame}{Bounding Instead of Point Identification}
If
\[
P(R=1\mid D=1)>P(R=1\mid D=0),
\]
trim the treated outcome distribution by the excess retention fraction.

This yields upper and lower bounds for the treatment effect among always-observed units.
\end{frame}

\begin{frame}{Why Bounds Matter}
When point identification is not credible, informative bounds may be more honest than a precise-looking adjusted estimate.
\end{frame}

\section{Principal-Stratification Intuition}

\begin{frame}{Potential Observation Status}
Define
\[
R_i(1),\qquad R_i(0).
\]

Units can be always observed, observed only if treated, observed only if untreated, or never observed.
\end{frame}

\begin{frame}{Always-Observed Effect}
One possible estimand is
\[
\E[Y(1)-Y(0)\mid R(1)=R(0)=1].
\]

Membership in the always-observed stratum is latent, so identification requires additional assumptions or bounds.
\end{frame}

\section{Selection Bias}

\begin{frame}{Sample Selection Before Treatment}
Sometimes units enter the analysis only if
\[
S=1.
\]

If selection depends jointly on treatment causes and outcome causes, conditioning on \(S=1\) can induce bias.
\end{frame}

\begin{frame}{Administrative Data Are Selected Data}
Examples:
\begin{itemize}
\item only customers who remain active;
\item only patients who return to clinic;
\item only workers who stay employed;
\item only firms that survive;
\item only users with trackable devices.
\end{itemize}
\end{frame}

\section{MNAR Sensitivity}

\begin{frame}{Pattern-Mixture Sensitivity}
A simple sensitivity model can posit
\[
\E[Y\mid R=0,D,X]
=
\E[Y\mid R=1,D,X]
+
\delta.
\]

Vary \(\delta\) over plausible values and recompute treatment effects.
\end{frame}

\begin{frame}{Selection-Model Sensitivity}
Alternatively model
\[
P(R=1\mid Y,D,X)
\]
directly, allowing missingness to depend on the unobserved outcome.
\end{frame}

\begin{frame}{Tipping-Point Analysis}
Ask:
\begin{itemize}
\item How different would missing outcomes need to be?
\item How differential would missingness need to be across treatment arms?
\item At what assumption value does the substantive conclusion change?
\end{itemize}
\end{frame}

\section{Applied Reasoning}

\begin{frame}{Example: Treatment Follow-Up}
Suppose a treatment reduces an outcome, but treated participants with poor outcomes are more likely to miss follow-up.

Observed treated outcomes become selectively optimistic.

A complete-case analysis can exaggerate treatment benefit.
\end{frame}

\begin{frame}{A Credible Missing-Data Workflow}
\begin{enumerate}
\item Define which variables are missing and when.
\item Plot missingness by treatment and baseline covariates.
\item State the target population and estimand.
\item Draw the selection/missingness causal structure.
\item Decide whether MCAR/MAR is defensible.
\item Use weighting, imputation, or augmented methods under MAR.
\item Consider bounds when point identification is weak.
\item Run MNAR sensitivity analysis.
\item Report attrition rates and effective sample sizes.
\end{enumerate}
\end{frame}

\begin{frame}{Common Failure Modes}
\begin{itemize}
\item dropping missing cases without discussion;
\item interpreting randomized treatment as protection against attrition bias;
\item conditioning on post-treatment causes of missingness;
\item using unstable observation weights;
\item imputing outcomes without treatment or missingness predictors;
\item treating MAR as testable from observed data alone.
\end{itemize}
\end{frame}

\begin{frame}{Synthesis: Missingness Changes the Design}
\begin{itemize}
\item Missingness is a selection mechanism, not merely a data-cleaning nuisance.
\item Randomization does not protect against treatment-dependent attrition.
\item Complete-case analysis can condition on a collider and alter the estimand.
\item Weighting and imputation require assumptions about observation.
\item Augmented estimators can combine observation and outcome models.
\item Lee bounds offer partial identification under monotone selection.
\item MNAR mechanisms require sensitivity analysis or stronger assumptions.
\end{itemize}
\end{frame}

\begin{frame}{Selected References}
\small
\begin{itemize}
\item Rubin, D. B. (1976). Inference and Missing Data. \textit{Biometrika}, 63(3), 581--592.
\item Little, R. J. A., and Rubin, D. B. (2019). \textit{Statistical Analysis with Missing Data}, 3rd ed. Wiley.
\item Robins, J. M., Rotnitzky, A., and Zhao, L. P. (1995). Analysis of Semiparametric Regression Models for Repeated Outcomes in the Presence of Missing Data. \textit{JASA}, 90(429), 106--121.
\item Lee, D. S. (2009). Training, Wages, and Sample Selection: Estimating Sharp Bounds on Treatment Effects. \textit{Review of Economic Studies}, 76(3), 1071--1102.
\item Molenberghs, G., and Kenward, M. G. (2007). \textit{Missing Data in Clinical Studies}. Wiley.
\end{itemize}
\end{frame}

\contactslide

\end{document}
